\chapter{From finite tops to the Haagerup boundary ring}
\label{ch:branching}

We now know that $U$ is proper and strongly rational. These two facts
still leave many possible representation categories. The next step
uses information already visible near the bottom of the lattice
theory: six retained currents, eighty complementary currents, and
the two eighty-one-dimensional glue spaces. A universal bound on
lowest module weights makes that small information exhaustive.

The chapter follows a chain of implications. Fourteen modules of a
known core bound all $U$-module tops. That bound fixes the branching
of $B$ and $H$. The extension algebra then limits the boundary category
to rank four or rank six. Character and indicator arguments select the
Haagerup fusion ring. The final tensor-category identification still
needs unitarity and classification, so we keep it separate from the
ring calculation.

\section{Lowest weights of the fourteen core modules}

The regular subalgebra $A_L\subset U$ has category
$\mathrm{Fib}\boxtimes SO(3)_{13}$. Label its fourteen simple modules
by $x_a=(1,a)$ and $y_a=(\tau,a)$ for $0\leq a\leq6$; the actual
even affine label is $2a$.

Their top spaces can be calculated from a small lattice minimization.
The six Potts modules in Example~\ref{ex:c-potts-six-modules} give,
for an integer charge $m$,
\[
 p_1(m)=\begin{cases}0&3\mid m,\\2/3&3\nmid m,\end{cases}
 \qquad
 p_\tau(m)=\begin{cases}2/5&3\mid m,\\1/15&3\nmid m.\end{cases}
\]
The actual lowest lattice--Potts weight in charge $j$ is
\begin{equation}\label{eq:c-T-top}
 h_{\mathcal T}(\epsilon,j)=
 \min_{m\equiv j\ (13)}
 \left(\frac{m^2}{156}+\frac{\mathbf1_{m\text{ odd}}}{4}
                                      +p_\epsilon(m)\right).
\end{equation}
The parity term is the left level-one affine contribution. Even $m$
gives a left singlet and odd $m$ a doublet. The right Cartan charge
is $m$.

For $-6\leq k\leq6$, the grade calculation of
Section~\ref{sec:c-pf-grade} gives the parafermion lowest weight
\begin{equation}\label{eq:c-PF-top}
 h_{\mathrm{PF}}(a,k)=\frac{a(a+1)}{15}-\frac{k^2}{13}
                                      +\max(0,|k|-a).
\end{equation}
The last term is important. When the required Cartan charge lies
outside the finite affine highest-weight representation, it first
appears at a positive affine grade. A twist calculation modulo one
would lose precisely this integer correction.

To find an $A_L$ top, minimize the sum of
\eqref{eq:c-T-top} and \eqref{eq:c-PF-top} over the extension
charges, with $k=-5j$ modulo thirteen. This is genuinely finite.
For $h_{\mathcal T}$ in each charge class, taking $|m|\leq6$
supplies a candidate at energy at most $14/13$, whereas
$|m|\geq13$ already exceeds that bound from its quadratic term
alone. The parafermion contribution is fixed within the charge
class. Hence $-12\leq m\leq12$ contains every minimizer.

\begin{example}[Two top weights]
For $x_1$, the neutral charge gives zero lattice--Potts energy and
parafermion energy $1\cdot2/15=2/15$. To prove that it is the
minimum, first consider $|k|\geq2$. The parafermion contribution is
\[
 \frac2{15}-\frac{k^2}{13}+|k|-1>\frac2{15}
 \qquad(2\leq |k|\leq6),
\]
as follows by checking the minimum of the concave quadratic on this
integer interval at its endpoints. The remaining $k=\pm1$ require
$j=\pm5$. Their lattice--Potts minimum is $14/13$, already larger
than $2/15$. The neutral candidate therefore wins, and its top is
one dimensional.

For $y_3$, choose $m=2$, hence $j=2$ and $k=3$. The two parts are
\[
 \frac4{156}+\frac1{15}=\frac6{65},\qquad
 \frac{12}{15}-\frac9{13}=\frac7{65}.
\]
Their sum is $1/5$. The conjugate charge $m=-2$ is also a minimizer,
giving a two-dimensional top. This multiplicity has a concrete charge
origin; it is not its quantum dimension.
\end{example}

\section{\texorpdfstring{Energy traces and the universal $17/10$ bound}{Energy traces and the universal 17/10 bound}}
\index{lowest-weight bound!seventeen tenths}

Let $E$ be the lowest space of any simple ordinary $U$-module, of
weight $h$ and dimension $D$. The compatible-top lemma says that
$E$ is a direct sum of complete $A_L$ top representations. Let $x_a$
and $y_a$ now also denote their nonnegative integer multiplicities.
The meaning is unambiguous when they occur in sums of dimensions.

On a core top let $L,P,Z$ be the traces of the left affine, Potts,
and PF$_{13}$ energy operators, and let $Q=\Tr H_R(0)^2$.
Define the additive bookkeeping quantity
\begin{equation}\label{eq:c-top-W}
 W=L+\frac{83}{6084}Q+\frac{25}{6}P+\frac{25}{26}Z.
\end{equation}
Why is the coefficient of $Q$ correct? In a complete right
$\mathfrak{sl}_2$ representation of highest weight $r$,
\[
 \frac{\Tr H_R(0)^2}{\dim}=\frac{r(r+2)}3,
 \qquad L^{A_{1,39}}_0=\frac{r(r+2)}{164}.
\]
Thus the Sugawara trace on the complete space $E$ equals $3Q/164$.
Multiplying by the
positive-stress coefficient $3403/4563$ gives $83Q/6084$.
An individual core top need not be a complete right
$\mathfrak{sl}_2$ representation. Its table entry $W$ is therefore
not asserted to be an ambient-energy trace on that core top, or an
energy eigenvalue times its dimension. Only after all core tops are
assembled into the actual $E$ does their sum satisfy
$W_{\rm total}=\Tr_E L_0=hD$. This distinction is why an isolated
row can have $W/\dim$ greater than the final universal bound.

The finite minimization gives the following complete table:
\[
\begin{array}{c|r|c|r@{\qquad}c|r|c|r}
 &\dim&h_{A_L}&4680W&&\dim&h_{A_L}&4680W\\\hline
 x_0&1&0&0&y_0&1&2/5&7800\\
 x_1&1&2/15&600&y_1&7&8/15&36960\\
 x_2&5&2/5&10440&y_2&15&4/5&110880\\
 x_3&9&4/5&56640&y_3&2&1/5&4080\\
 x_4&2&1/3&5520&y_4&14&11/15&81480\\
 x_5&14&1&148560&y_5&6&2/5&18240\\
 x_6&6&4/5&54240&y_6&2&1/5&4920
\end{array}
\]
This table is a finite arithmetic input obtained from the actual
core realization. It can be reconstructed from the two minimization
formulas and charge multiplicities; it does not assert any list of
$U$-modules in advance.

It is instructive to recover the charges that are hidden by the
three numerical columns. In the following table, an entry lists all
minimizing absolute charge classes $|j|$; both signs are included
when $j\ne0$:
\[
\begin{array}{c|l@{\qquad}c|l}
x_0&0&y_0&0\\
x_1&0&y_1&0,5\\
x_2&0,3&y_2&0,2,3,5\\
x_3&0,2,3,6&y_3&2\\
x_4&6&y_4&1,2,4,6\\
x_5&1,4,6&y_5&1,4\\
x_6&4&y_6&4
\end{array}
\]
One must keep tied representatives of a charge class. For
$\epsilon=1$ the ties at $j=1,4,5$ are respectively
$m=1,-12$, $m=4,-9$, and $m=5,-8$. For $\epsilon=\tau$,
the ties at $j=3,5,6$ are $m=3,-10$, $m=5,-8$, and $m=6,-7$.
Negative $j$ gives the opposite charges. An odd $m$ contributes
a left doublet, so it contributes two states at that right charge.

For example, the $j=\pm1$ ties of $x_5$ supply charges
$\pm1$ with multiplicity two and $\pm12$ with multiplicity one.
Its $j=\pm4$ ties supply charges $\pm4$ with multiplicity one
and $\pm9$ with multiplicity two. The $j=\pm6$ contribution
supplies one state at each of $\pm6$. The total is
$2(2+1+1+2+1)=14$, reproducing the dimension column.

Completeness of the right current representations imposes two further
inequalities,
\begin{equation}\label{eq:c-top-inequalities}
 y_2\geq x_5,\qquad y_4\geq x_5+x_6.
\end{equation}
The minimizers just listed give, for the whole space $E$,
\[
 \operatorname{mult}(12)=x_5,\quad
 \operatorname{mult}(10)=y_2,\quad
 \operatorname{mult}(9)=2(x_5+x_6),\quad
 \operatorname{mult}(7)=2y_4.
\]
In a finite $\mathfrak{sl}_2$ representation the multiplicity of
a nonnegative weight cannot increase as the weight increases by
two. Comparing $10$ with $12$ and $7$ with $9$ gives exactly
\eqref{eq:c-top-inequalities}. The factors of two in the odd
charges are the left doublets. A right-current representation cannot
be assembled from isolated Cartan eigenlines while omitting its
required lower-weight companions.

The following nonnegative certificate packages the entire bound:
\begin{align}\label{eq:c-top-certificate}
4680\left(\frac{17}{10}D-W_{\rm total}\right)
={}&8460(y_2-x_5)+28716(y_4-x_5-x_6)\notag\\
 &+7956x_0+7356x_1+29340x_2+14964x_3+10392x_4\notag\\
 &+22212x_6+156y_0+18732y_1+11832y_3\notag\\
 &+1188y_4+29496y_5+10992y_6\geq0.
\end{align}
Each coefficient can be checked directly from the table. For example,
the $y_0$ coefficient is $4680(17/10)-7800=156$.
For $x_5$, the right side contributes
$-8460-28716=-37176$, agreeing with
$4680(17/10)\cdot14-148560$.
This is a small linear-programming dual certificate written as an
identity. The two representation constraints absorb the troublesome
negative coefficient, and every remaining term is visibly
nonnegative. No optimization program is needed to verify it.

\begin{proposition}[Every simple top lies below weight two]
Every simple ordinary $U$-module has lowest weight
\[
 0\leq h\leq\frac{17}{10}<2.
\]
\end{proposition}
\begin{proof}
The preceding chapter gives $h\geq0$. Trace the positive-stress
identity on $E$: compact averaging preserves trace, so
$W_{\rm total}=hD$. Equations \eqref{eq:c-top-inequalities} and
\eqref{eq:c-top-certificate} give the upper bound.
\end{proof}

The bound need not be optimal. Its decisive feature is being less
than two. That excludes a new integer-graded summand that first
appears at weight two or later.

\section{\texorpdfstring{Why eighty currents determine the whole $B$ branching}{Why eighty currents determine the whole B branching}}
\index{branching!lattice extension}

Since $U$ is strongly rational, its restriction inside $B$ is
semisimple. Every summand has integer lowest weight, because $B$ is
integer graded. The universal bound therefore permits only weights
zero and one. The weight-zero summand is the vacuum, occurring once.

There are $86-6=80$ remaining currents in $B_1$. The finite
construction provides an irreducibility certificate for their
zero-mode action. Its logic is elementary and valuable.

\begin{lemma}[A cyclic one-dimensional eigenline]
Let a finite-dimensional positive Hermitian space carry a
star-closed algebra of operators containing commuting self-adjoint
Cartan operators. If one joint Cartan eigenspace is one dimensional
and a nonzero vector on that line is cyclic, the representation is
irreducible.
\end{lemma}
\begin{proof}
An invariant subspace has an invariant orthogonal complement because
the algebra is star closed. Both are Cartan invariant. The
one-dimensional joint eigenspace must lie entirely in one of them.
Cyclicity makes that summand the entire space.
\end{proof}

\begin{constructioninput}[Five irreducible finite tops]
On the eighty-dimensional current complement, the known $U$ zero
modes are star closed and a vector in the one-dimensional Cartan
weight $(0,10)$ line is cyclic. In each of the four mixed nontrivial
$E_6\oplus A_2$ lattice sectors of lowest weight one, the analogous
star-closure, one-dimensional eigenline, and cyclicity tests prove
irreducibility of its eighty-one-dimensional top.
\end{constructioninput}

There are four mixed sectors because both $\Z/3$ discriminant
coordinates are nonzero. The $E_8$ extension chooses two of these
$B$-module sectors as its glue sectors.
The bound rules out any summand invisible in these tops. Hence
\begin{equation}\label{eq:c-B-branching}
 B\downarrow_U=U\oplus X,
 \qquad h_X=1,\quad\dim X_1=80.
\end{equation}
The two actual glue sectors in $H/B$ restrict to simple modules $Y,Y^*$
with $81$-dimensional tops, giving
\begin{equation}\label{eq:c-H-branching}
 H\downarrow_U=\one\oplus X\oplus Y\oplus Y^*.
\end{equation}
The notation $Y^*$ means the contragredient. It does not yet say
whether $Y$ and $Y^*$ are isomorphic.

\section{The extension algebra remembers a boundary category}

Let $\mathcal C=\Rep(U)$ and regard $H$ as its commutative extension
algebra. Because $H$ is holomorphic, this algebra is Lagrangian:
its local-module category is trivial and its dimension is the square
root of the global dimension of $\mathcal C$.
Let
\[
 \mathcal F=\mathcal C_H
\]
be the category of \emph{all} right $H$-modules. Local modules and
all modules must not be confused: the former are trivial here, while
the latter will be the desired nontrivial boundary category.
Theorem~\ref{thm:lagrangian-reconstruction} gives
$\mathcal C\simeq Z(\mathcal F)$.

Temporarily write $\sigma=\FPdim_{\mathcal C}B$. Extension
dimension identities give
\[
 \qdim X=\sigma-1,\qquad\qdim Y=\sigma,
 \qquad\operatorname{Dim}\mathcal C=9\sigma^2,
 \quad\operatorname{Dim}\mathcal F=3\sigma.
\]
Here $\operatorname{Dim}$ means global dimension, the sum of squared
simple quantum dimensions. Positivity of the nonvacuum conformal
weights gives the positive spherical structure through character
asymptotics. At this stage pseudo-unitarity is sufficient; a
categorical $C^*$ structure is supplied later.

For center induction $I$ and forgetful functor $F$, the canonical
algebra is $I(\one)$. The induction identity
$FI(\one)=\bigoplus_z z\otimes z^*$, over the simple boundary
objects $z$, implies
\[
 \operatorname{rank}\mathcal F
 =\dim\End_{\mathcal C}(H)
 =\begin{cases}4&Y\not\cong Y^*,\\6&Y\cong Y^*.\end{cases}
\]
Indeed, the endomorphism dimension of a semisimple object is the sum
of the squares of its multiplicities. Four distinct summands give
$1+1+1+1=4$; identifying the two glue summands gives $1+1+2^2=6$.

Why should an automorphism of an extension say anything about a
boundary particle? The canonical-algebra theorem gives a useful
answer that does not yet require a unitary tensor category.

\begin{lemma}[Extension automorphisms and invertible boundary objects]
\label{lem:c-auto-invertible}
For a pseudo-unitary fusion category $\mathcal F$, the algebra
automorphism group of $I(\one)\in Z(\mathcal F)$ injects into
the group of invertible objects of $\mathcal F$.
\end{lemma}
\begin{proof}
The canonical algebra has a normalized special Frobenius structure
with multiplication $m$, comultiplication $\Delta$, and
$m\Delta=\mathrm{id}$. Besides composition, its endomorphisms carry
the convolution product
$f*g=m(f\otimes g)\Delta$.
The standard two-product description identifies composition with
multiplication in $K_0(\mathcal F)\otimes\C$, while the primitive
convolution idempotents are $d_z[z]$, with $d_z>0$ the simple
dimensions.

An algebra automorphism obeys $f*f=f$, and its composition inverse
does too. Consequently their fusion-algebra expansions have
coefficients either zero or $d_z$. Their fusion product is
$[\one]$. Nonnegative fusion coefficients prevent cancellation,
so every pair of objects in their supports must multiply only to
copies of the unit. Rigidity implies that such a pair consists of
inverse invertible simples. Both supports are therefore singletons,
with dimension one. Composition respects the assignment, proving
the group injection.
\end{proof}

The actual order-three automorphisms of $H/B$ therefore supply three
invertible objects in $\mathcal F$. More is known than their fusion group:
their associator is untwisted. The discriminant form of $B$ is
\[
 q(a,b)=\frac{2a^2+b^2}{3},\qquad(a,b)\in(\Z/3)^2.
\]
Its diagonal and antidiagonal isotropic lines are complementary.
Induction from the latter across the extension selected by the former
gives a full tensor copy of $\mathrm{Vec}_{\Z/3}$ in $\mathcal F$.
Knowing only that there are three invertible objects would not rule
out a nontrivial associator class.

\section{Excluding a sign extension and the rank-four fork}
\index{rank-four obstruction}

First, $X$ cannot be invertible. Otherwise the extension $B=U\oplus X$
would have an order-two grading and a sign automorphism $g$ with
$B^g=U$. Its fixed current Lie algebra would be
$\mathfrak{sl}_2\oplus\mathfrak{sl}_2$, of rank two.
But every finite-order automorphism of $E_6$ fixes a toral subalgebra
of dimension at least four, and every such automorphism of $A_2$
fixes one of dimension at least one. The nonisomorphic ideals cannot
be exchanged. The fixed rank is therefore at least five, a
contradiction. This also excludes the possibility that all six
boundary simples are invertible.

At rank four there are three invertibles and one noninvertible, with
near-group ring
\[
 \alpha^3=1,\quad\alpha\rho=\rho\alpha=\rho,
 \qquad\rho^2=1+\alpha+\alpha^2+m\rho.
\]
The relevant pseudo-unitary near-group theorem restricts $m$ to
$0,2,3,6$. This is a specialized categorical theorem, not a consequence
of associativity of the displayed ring alone.

The case $m=0$ makes the extension $B/U$ have dimension two, so $X$
would be invertible. For $m=2$, the categorical reduction gives a
possibly twisted $A_4$ double and a holomorphic charge-eight extension
whose $A_4$ fixed algebra is $U$. The holomorphic extension is $E_8$.
An elementary current estimate then excludes it. The adjoint traces
of nontrivial order-two and order-three automorphisms of $E_8$ are
bounded below by $-8$ and $-4$, respectively. Since $A_4$ has three
involutions and eight elements of order three,
\[
 \dim(E_8^{A_4})\geq
 \frac{248+3(-8)+8(-4)}{12}=16>6.
\]

The remaining two cases are excluded by the following specialized
lemma from the construction's modular-character argument. We state
its hypotheses because the conclusion is not a general restriction
on arbitrary near-group categories.

\begin{lemma}[Rank-four character obstruction]\label{lem:c-rank4-character}
Suppose a strongly rational charge-eight VOA has positive nonvacuum
lowest weights, an extension $B=V_{E_6\oplus A_2}$ with branching
$B=U\oplus X$, and the rank-four pseudo-unitary near-group center
obtained above, with its untwisted order-three subgroup and the
specified lattice extension. For multiplicities $m=3,6$, its current
dimension must have one of the following forms:
\[
\begin{array}{c|l}
 m&\text{necessary number of currents}\\\hline
 3&15+A+10B\quad\text{or}\quad14+42A\\
 6&8+A\quad\text{or}\quad5+12P
\end{array}
\qquad A,B,P\in\Z_{\geq0}.
\]
The letters $A,B,P$ in this table are weighted dimensions of polar
module tops, not the previously named vertex algebras.
\end{lemma}

The integers are nonnegative because they count actual top states.
None of the displayed forms equals six. We explain the mechanism of
the specialized lemma next; the assertion includes the finite
theta-coefficient calculation rather than treating its entries as
new lattice construction assumptions.

The mechanism deserves explanation even without reproducing a theta
atlas. A \emph{polar} character term is a negative power of $q$,
equivalently a state of weight below $c/24$; the cusp is $q=0$,
or $\tau\to i\infty$. Center induction determines a finite modular
representation
for a suitable difference of branching characters. Multiplication by
$\eta^8$ makes it a weight-four theta-type function. The allowed polar
data, before that multiplication, are specified by the nonnegative
module-top dimensions above. Once those are prescribed, principal-part cancellation determines the
function; its vacuum coefficient yields the displayed restrictions.
One keeps the factor of two between the vacuum character and the
sum/difference branching coordinates. Dropping it changes the current
obstruction. This is an application of modular functions to a finite
categorical alternative, not an assumption that the desired character
vector was known from the start.

\subsection{A complete theta cancellation in one rank-four alternative}

We work the first $m=3$ orientation in detail. This demonstrates
where the integers $15$ and $10$ come from, and lets us check
the potentially dangerous vacuum factor of two.

The categorical projection in this orientation is indexed like the
six integrable $G_{2,3}$ modules, whose finite-dimensional labels
are $1,7,27,77,14,64$, tensored with the three $A_{2,1}$
labels. For each pointed $A_2$ charge, retain the four-component
difference vector
\[
 v=(\chi_1-\chi_{64},-\chi_{77},-\chi_{14},-\chi_7).
\]
Here the labels describe the modular projection fixed by the
rank-four induction argument; they do not assert that the unknown
VOA is already an affine algebra. The known lattice extension
gives $\chi_1+\chi_{64}=\chi_B$ in the pointed vacuum charge.
Since $\chi_1=\chi_U$, a difference between two solutions of
this extension problem obeys
\[
 \delta v_{0,0}=2\,\delta\chi_U.
\]
The lattice character is fixed, so the plus and minus coordinates
cannot vary independently.

After multiplication by $\eta^8$, the projected modular
representation is the even Weil representation on
$(r,s)\in(\Z/7)\times(\Z/3)$, even in each coordinate.
Representatives are $0\leq r\leq3$ and $s=0,1$. Its $T$
phases are $\exp(2\pi i(r^2/7+s^2/3))$. One convenient
normalization of $S$ is the tensor product of
\[
\begin{gathered}
 S^{(7)}_{00}=-1/\sqrt7,\qquad
 S^{(7)}_{0r}=-2/\sqrt7,\qquad
 S^{(7)}_{r0}=-1/\sqrt7,\\
 S^{(7)}_{rs}=-2\cos(4\pi rs/7)/\sqrt7
 \quad(r,s>0)
\end{gathered}
\]
with the even order-three Fourier matrix
$3^{-1/2}\left(\begin{smallmatrix}1&2\\1&-1\end{smallmatrix}\right)$.
The unbalanced first row and column simply record that a nonzero
even coordinate represents two finite-group elements.

We construct weight-four forms for this representation using the
rank-four lattice $L_7\oplus A_2$, where
\[
 \operatorname{Gram}(L_7)=
 \begin{pmatrix}2&1\\1&4\end{pmatrix}.
\]
Let $u,v$ be its displayed basis and put $w=v-u/2$, so
$u\perp w$, $u^2=2$, $w^2=7/2$. Write a vector in its
real plane as $pu+tw$, and put
\[
 a=p^2,\qquad b=7t^2/4,\qquad c=(y,y)/2
 \quad(y\in A_2\otimes\R).
\]
These lowercase letters are temporary quadratic coordinates.
Both $H_1=a-b$ and $H_2=a+b-c$ are harmonic homogeneous
quadratics. The two coordinates in each plane contribute equal
Laplacians, which cancel in these differences. Poisson summation
therefore makes the coset theta sums
\[
 \Theta_{H,\lambda}(\tau)
    =\sum_{x\in L+\lambda}H(x)q^{(x,x)/2}
\]
vector-valued modular forms of weight $4/2+2=4$.

Only three small coefficients are needed:
\[
\begin{array}{c|c|r|r}
 (r,s)&q\text{-power}&\Theta_{H_1}&\Theta_{H_2}\\\hline
 (1,0)&1/7&-1/7&1/7\\
 (3,0)&2/7&3/7&4/7\\
 (0,0)&1&2&-4
\end{array}
\]
Each entry can be counted by hand. The first coset has shortest
vector $2w/7$, giving $a=0,b=1/7$. In the second there are
two shortest vectors with $a=1/4,b=1/28$, each of energy
$2/7$; hence their total $H_1$ value is $3/7$ and their
total $H_2$ value is $4/7$. In the vacuum coset, the two
$L_7$ roots contribute $2$ to both forms. The six $A_2$
roots contribute zero to $H_1$ and $-6$ to $H_2$.

The only allowed nonvacuum polar top dimensions are $A$ at
weight $2/7$ and $B$ at weight $1/7$. Either can be zero
if the actual module begins an integer grade higher. Use
$G_{2,3}\otimes A_{2,1}$ as a reference solution of the same
modular projection and lattice-extension problem. Its central
charge is $3\cdot14/(3+4)+2=8$, its currents number
$14+8=22$, and its polar dimensions are $A=7,B=0$.
The fundamental seven-dimensional $G_2$ representation has
quadratic Casimir two in this normalization, so its Sugawara
weight is $2/(3+4)=2/7$.
Thus the polar coefficients of $\eta^8\delta v$ are
$7-A$ at $(3,0)$ and $-B$ at $(1,0)$.

Matching them against
$\lambda\Theta_{H_1}+\mu\Theta_{H_2}$ gives two linear
equations,
\[
 \frac{-\lambda+\mu}{7}=-B,\qquad
 \frac{3\lambda+4\mu}{7}=7-A.
\]
Their solution is
$\lambda=7-A+4B$, $\mu=7-A-3B$.
After division by $\eta^8$, the unmatched remainder has no
poles at the cusp. The finite-image modular representation
makes it a holomorphic vector of modular functions on a compact
finite cover, so it is constant. Its only possible $T$-fixed
coordinate is $(r,s)=(0,1)$, and the displayed Fourier $S$
matrix does not preserve that line. Hence the constant is zero:
the matched theta expression is the whole difference, not merely
its first few coefficients.

Finally the vacuum coefficient of the theta expression is
$2\lambda-4\mu$. By the vacuum factor of two above it is also
$2(\dim U_1-22)$; the leading vacuum terms canceled before
multiplication by $\eta^8$, so no further eta coefficient
enters this difference. We obtain
\[
 2(\dim U_1-22)
 =2(7-A+4B)-4(7-A-3B),
 \qquad \dim U_1=15+A+10B.
\]
Since $A,B$ are nonnegative integers, six currents are impossible.
The other three alternatives in Lemma~\ref{lem:c-rank4-character}
are the corresponding specialized theta cancellations for the reverse
orientation and the $m=6$ projected representation. The complete
calculation here explains the method and fixes its normalization;
the remaining coefficient formulas are retained as the stated
specialized lemma.

\section{Deriving the six-object fusion ring}

Rank four is excluded, so
$H=\one\oplus X\oplus2Y$. Its endomorphism algebra is
\[
 \End_{\mathcal C}(H)\cong\C\oplus\C\oplus M_2(\C).
\]
The two-product description of the canonical center algebra identifies
this composition algebra with the complexified boundary fusion algebra.
In particular the fusion ring is noncommutative.

The group of three invertibles acts on the three noninvertibles. An
order-three permutation is either trivial or transitive. The trivial
case would split off the two nontrivial characters of $\C[\Z/3]$
as one-dimensional ideals, leaving an $M_2(\C)$ ideal on which the
dimension character is nonzero. A matrix algebra has no nonzero
one-dimensional representation, so this is impossible.

Choose a noninvertible $\rho$ and label the orbit by
$\rho,\alpha\rho,\alpha^2\rho$. Duality conjugates the left and
right actions of the invertible group, so the right action is
transitive too. It commutes with the left action and hence is the
same three-cycle or its inverse; the identity permutation has now
been excluded. The same cycle
would make the algebra generated by $\alpha,\rho$ commutative. Thus
$\rho\alpha=\alpha^2\rho$. To see the self-duality explicitly, write
$\rho^*=\alpha^k\rho$. Then
$(\alpha^i\rho)^*=\alpha^{i+k}\rho$; applying duality twice gives
$2k=0$ modulo three. Hence $k=0$, and all three noninvertibles
are self-dual. Frobenius reciprocity gives exactly one invertible in
$\rho^2$, namely the unit; commuting with $\alpha$ equates the remaining
three coefficients. Therefore
\begin{equation}\label{eq:c-HI-m}
 \alpha^3=1,\quad\rho\alpha=\alpha^2\rho,
 \qquad\rho^2=1+m(\rho+\alpha\rho+\alpha^2\rho),\quad m\geq1.
\end{equation}
There is still a multiplicity $m$ to determine. Recognizing the shape
of the Haagerup ring is not enough to set it to one.

\section{Indicators: rotating a junction without braiding}
\label{sec:c-indicator-bridge}
\index{Frobenius--Schur indicators!definition}

Let $\mathcal C$ be a spherical fusion category over $\C$, with a
fixed pivotal identification $j_X:X\to X^{**}$. A unitary fusion
category has a canonical positive spherical structure. The theorem
used here needs the spherical structure, without assuming the
complete unitarity proved later in Chapter~\ref{ch:unitarity}.
For $n\ge1$, consider the finite-dimensional junction space
$\Hom(\one,X^{\otimes n})$. Bend its first outgoing strand around
the junction and return it as the last strand, using evaluation,
coevaluation, and $j_X^{-1}$. With the required reassociations this
defines the \emph{pivotal cyclic rotation} $R_{X,n}$.
It uses no crossing and requires no braiding on $\mathcal C$.
In formulas, put $Y=X^{\otimes(n-1)}$ and let
$f:\one\to X\otimes Y$. The two stages are
\begin{align*}
 \widetilde f&=(\operatorname{ev}_X\otimes1_Y)
                    (1_{X^*}\otimes f):X^*\to Y,
 \\
 R_{X,n}(f)&=(\widetilde f\otimes j_X^{-1})
                    \operatorname{coev}_{X^*}:\one\to Y\otimes X,
\end{align*}
with the same associators understood as in the cup and cap identities.
The $n$th Frobenius--Schur indicator is the ordinary linear trace
\begin{equation}\label{eq:c-indicator-definition}
 \nu_n(X)=\Tr_{\Hom(\one,X^{\otimes n})}(R_{X,n}).
\end{equation}
Pivotal coherence gives $R_{X,n}^{\,n}=1$. For simple $X$, this
immediately explains two values used below:
\[
 \nu_1(X)=\delta_{X,\one},\qquad
 \nu_2(X)=
 \begin{cases}
 0,&X\not\simeq X^*,\\
 +1\text{ or }-1,&X\simeq X^*.
 \end{cases}
\]
Indeed, $\Hom(\one,X^{\otimes2})\simeq\Hom(X^*,X)$ is zero or
one-dimensional, and its rotation squares to the identity. The sign
measures the symmetry of the chosen self-duality under pivotal
transpose; it is not fixed merely by the equation $X\simeq X^*$.

\begin{example}[The group formula from an invariant projector]
Take $\mathcal C=\Rep(G)$ for a finite group, with its usual tensor
product and positive spherical structure. For a representation
$D:G\to\operatorname{GL}(V)$, the junction space is
$(V^{\otimes n})^G$. Pivotal rotation is the cyclic permutation
\[
 C(v_1\otimes\cdots\otimes v_n)
   =v_2\otimes\cdots\otimes v_n\otimes v_1.
\]
The projector onto invariants is
$P=|G|^{-1}\sum_{g\in G}D(g)^{\otimes n}$, and $C$ commutes with
$P$. Therefore its trace on invariants equals $\Tr(CP)$ on the
whole tensor power. For any matrix $A$, expansion in a tensor basis
gives
\[
 \Tr(CA^{\otimes n})
 =\sum_{i_1,\ldots,i_n}
    A_{i_1i_2}A_{i_2i_3}\cdots A_{i_ni_1}
 =\Tr(A^n).
\]
Thus, writing $\chi_V(g)=\Tr D(g)$, we obtain the classical formula
\begin{equation}\label{eq:c-group-indicator}
 \nu_n(V)=\frac1{|G|}\sum_{g\in G}\chi_V(g^n).
\end{equation}
For $G=\Z_3=\langle t\rangle$ and the one-dimensional character
$\chi(t)=\zeta_3=e^{2\pi i/3}$, this becomes
\[
 \nu_n(\chi)=\frac{1+\zeta_3^n+\zeta_3^{2n}}3
 =\begin{cases}1,&3\mid n,\\0,&3\nmid n.\end{cases}
\]
In particular $\nu_3(\chi)=1$, although $\nu_1(\chi)=\nu_2(\chi)=0$.
The invariant projector and the cyclic trace identity explain this
without any modular data.
\end{example}

\begin{example}[A unitary self-dual object with negative second indicator]
The quaternion group $Q_8=\{1,-1,\pm i,\pm j,\pm k\}$ has a
unitary irreducible representation $V$ of dimension two, with character
$\chi_V(1)=2$, $\chi_V(-1)=-2$, and zero on the other six elements.
It is self-dual: its character is real. The two central elements
square to $1$, while the other six square to $-1$. Thus
\[
 \nu_2(V)=\frac{2\cdot2+6\cdot(-2)}8=-1.
\]
The invariant bilinear pairing is alternating rather than symmetric.
Its sign does not conflict with the positive Hermitian inner product
that makes the representation unitary; these are different pairings.
Consequently even a positive unitary setting does not allow us to
replace the sign $\varepsilon=\nu_2(\rho)$ below by $+1$ without
an argument.
\end{example}

The bridge to modular data is the following structural theorem.
Let $F:Z(\mathcal C)\to\mathcal C$ be forgetful and $I$ its
two-sided adjoint. Equip the center with the ribbon twist induced
by the same spherical structure. Then, for every object $X$ and
positive integer $n$,
\begin{equation}\label{eq:c-indicator-center}
 \nu_n(X)
 =\frac{\Tr_{\mathrm{cat}}(\theta_{I(X)}^{\,n})}
        {\operatorname{Dim}(\mathcal C)},\qquad
 \operatorname{Dim}(\mathcal C)=\sum_{a\text{ simple}}d_a^2.
\end{equation}
This is the center-induction theorem of Ng--Schauenburg. Here the
numerator is a \emph{categorical} trace: if
$I(X)=\bigoplus_Z b_ZZ$, then it is
$\sum_Z b_Zd_Z\theta_Z^n$. The twist is $\theta_Z$, without the
common central-charge phase of a VOA's modular $T$ matrix.

\begin{example}[Checking the normalization in the pointed center]
For $\mathcal C=\operatorname{Vec}_{\Z_3}$ with trivial associator
and positive spherical structure, let $\alpha$ be the simple object
of degree one. The center simples from Chapter~\ref{ch:centers}
are $(g,k)$, with $d_{(g,k)}=1$, twist $\zeta_3^{gk}$, and
$F(g,k)=g$. Adjunction gives
\[
 I(\alpha)=(1,0)\oplus(1,1)\oplus(1,2),\qquad
 \frac{\Tr_{\mathrm{cat}}(\theta_{I(\alpha)}^{\,n})}{3}
 =\frac{1+\zeta_3^n+\zeta_3^{2n}}3.
\]
Directly, $\Hom(\one,\alpha^{\otimes n})$ is zero unless $3\mid n$;
when nonzero it is one-dimensional and the untwisted pivotal rotation
is the identity. This verifies both sides of
\eqref{eq:c-indicator-center}. The fusion relation $\alpha^3=\one$
alone would not prove $\nu_3(\alpha)=1$: a nontrivial associator
can give a nontrivial rotation phase on that same one-dimensional
junction space.
\end{example}

For the rank-six fusion ring just derived, let $d_m$ be the positive
solution of $d_m^2=1+3md_m$. The simple dimensions are
$1,1,1,d_m,d_m,d_m$. Put $s_m=1+d_m^2$. Its global dimension is
$3s_m$, so the theorem says exactly
\[
 \Tr_{\mathrm{cat}}(\theta_{I(\rho)}^{\,n})
   =3s_m\nu_n(\rho).
\]
The first moment is therefore zero, and the second is
$3s_m\varepsilon$ with $\varepsilon=\nu_2(\rho)\in\{1,-1\}$.
There is no extra dimension factor $d_m$ on the right. On the left,
a residual induced summand $Z_j$ of multiplicity $n_j$, dimension
$3d_mn_j$, and twist $\theta_j$ contributes $3d_mn_j^2\theta_j^n$.
This is the source
of both the normalization and the squared multiplicities in the
following argument.

\section{Twist moments force multiplicity one}
\index{Frobenius--Schur indicators!twist moments}

Put $d_m=(3m+\sqrt{9m^2+4})/2$, so
$d_m^2=1+3md_m$, and put $s_m=1+d_m^2$.
We use the positive spherical structure already obtained above from
character asymptotics. No categorical adjoint is needed for the
following positive-dimension trace argument.

First derive the part of center induction that contains invertibles.
Write $R=\rho+\alpha\rho+\alpha^2\rho$. The underlying-object formula
$FI(x)=\bigoplus_z z\otimes x\otimes z^*$ and the fusion ring give
\[
 FI(\alpha)=FI(\alpha^2)=3\alpha+3\alpha^2+3mR.
\]
For example, the three invertible $z$ contribute $3\alpha$ to
$FI(\alpha)$; the three self-dual noninvertibles each contribute
$\alpha^2+mR$. By adjunction,
\[
 \dim\End I(\alpha)=\dim\End I(\alpha^2)
 =\dim\Hom(I(\alpha),I(\alpha^2))=3.
\]
An endomorphism dimension of three is a sum of squared nonnegative
integral multiplicities, so each induction object is a multiplicity-free
sum of three simples. Their mutual inner product three makes them
the same sum, say $C_0+C_1+C_2$. They are disjoint from $I(\one)$.
Adjunction gives one copy of each $\alpha,\alpha^2$ in $F(C_j)$.
Centrality makes its three noninvertible multiplicities equal. Hence
\[
 F(C_j)=\alpha+\alpha^2+b_jR,\qquad
 \sum_jb_j=3m,\qquad d_{C_j}=2+3b_jd_m.
\]

The pointed subcategory has trivial associator in this construction,
so $\nu_3(\alpha)=1$. The third twist trace therefore says
\[
 \sum_j d_{C_j}\theta_{C_j}^3
   =3s_m=\sum_jd_{C_j}.
\]
The twists have modulus one and the dimensions are positive. Equality
in the triangle inequality forces each $\theta_{C_j}^3=1$.
This does not yet say that all three cubic roots occur.
The additional identity $\nu_1(\alpha)=0$ gives
\[
 2\sum_j\theta_{C_j}+3d_m\sum_jb_j\theta_{C_j}=0.
\]
The number $d_m$ is a real quadratic irrational, while the twists now
lie in $\Q(\zeta_3)$. These two quadratic fields intersect in $\Q$,
so both sums vanish separately. Three cubic roots summing to zero
must be $1,\zeta_3,\zeta_3^2$, once each. The weighted sum then
forces $b_0=b_1=b_2=m$. Thus all three $C_j$ have dimension $s_m$.

Together with $I(\one)=\one+X+2Y$, this supplies six central
simples of dimensions $1,d_m^2$ and four copies of $s_m$, with
twists $1,1,1,1,\zeta_3,\zeta_3^2$. All central summands whose
forgetful object contains an invertible have now been exhausted.
Every remaining central simple $Z_j$ has twist $\theta_j$ and forgets
to $n_jR$, with $n_j$ a positive integer; centrality again equates its
three
noninvertible multiplicities.
The global-dimension identity and the first two Frobenius--Schur
indicator identities give
\begin{align}
 \sum_jn_j^2&=4m^2+2,\label{eq:c-indicator-count}\\
 \sum_jn_j^2\theta_j&=-\frac{m^2+m\sqrt{9m^2+4}}2,
 \label{eq:c-indicator-first}\\
 \sum_jn_j^2\theta_j^2&=\frac{-m^2+(2\varepsilon-m)
                            \sqrt{9m^2+4}}2,
 \qquad\varepsilon\in\{1,-1\}.
 \label{eq:c-indicator-second}
\end{align}
The roots of unity in these sums are counted with multiplicity
$n_j^2$. Distinctness is not assumed.

Here is how the normalizations enter. Denote the known nontrivial
summands of $I(\one)$ by $X,Y$, of dimensions $d_m^2,s_m$, and
the three summands of $I(\alpha)$ by $C_0,C_1,C_2$, each of
dimension $s_m$. Their twists are respectively $1,1$ and
$1,\zeta_3,\zeta_3^2$, as just derived using both the first and third
indicators. The center's global dimension is the square
of $3s_m$. Therefore
\[
 9s_m^2=1+d_m^4+4s_m^2+9d_m^2\sum_jn_j^2.
\]
Using $d_m-d_m^{-1}=3m$ reduces this to
$\sum_jn_j^2=4m^2+2$.

Adjunction gives the induction decomposition
\[
 I(\rho)=mX+mY+m(C_0+C_1+C_2)+\sum_jn_jZ_j.
\]
The first indicator of a nonunit simple is zero. Thus the first
twist trace of this object is
\[
 0=m(d_m^2+s_m)+3d_m\sum_jn_j^2\theta_j,
\]
because the three $C_j$ contributions cancel. This yields
$-(m/3)(2d_m+d_m^{-1})$ for the sum, which is the first moment
in \eqref{eq:c-indicator-first}. The second twist trace equals
$3s_m\varepsilon$, where $\varepsilon=\nu_2(\rho)$ is $1$ or
$-1$ for a self-dual simple. Again the $C_j$ contributions cancel:
\[
 3s_m\varepsilon
   =m(d_m^2+s_m)+3d_m\sum_jn_j^2\theta_j^2.
\]
Subtract the first equation and use
$s_m/d_m=d_m+d_m^{-1}=\sqrt{9m^2+4}$ to obtain
\eqref{eq:c-indicator-second}.

Here is the arithmetic force of the equations. Write
$9m^2+4=t^2D$ with $D$ squarefree. The allowed squareclasses have
strong restrictions. Since
$(3m)^2<9m^2+4<(3m+1)^2$, one has $D>1$.
Reduction modulo three shows that $3\nmid t$ and $D\equiv1\pmod3$.
Also $7\nmid D$: divisibility by seven would require
$m^2\equiv5\pmod7$, which is impossible.

Here $\varphi(n)$ is Euler's totient, the number of residues modulo $n$
that are coprime to $n$.
Larson's root-sum bound is a theorem about expressing a quadratic
irrationality as a sum of roots of unity; it is stronger than the
triangle inequality. Applying it to two copies of the root multiset
clears the half-integral coefficients in
\eqref{eq:c-indicator-first} and gives
\[
 4m^2+2\geq\frac{mt}{2}\varphi(2D).
\]
For precisely these squareclasses, except $D=10$, elementary
prime-factor estimates give
$\varphi(2D)/\sqrt D\geq12/\sqrt{13}$.
For odd $D$, this ratio is the product of
$(p-1)/\sqrt p$ over its distinct prime factors. A single allowed
prime is at least thirteen; two distinct prime factors contribute
at least $40/\sqrt{55}$, from five and eleven, which already
exceeds $12/\sqrt{13}$. Every additional factor is greater than one.
For even $D$, an additional factor $\sqrt2$ occurs. The smallest
single odd prime is five, giving the exception $D=10$; otherwise
it is at least eleven, and two odd factors give an even larger bound.
These restrictions matter: the inequality is false for arbitrary
squarefree numbers such as two or five. Consequently
\[
 4m^2+2\geq\frac{6m\sqrt{9m^2+4}}{\sqrt{13}}.
\]
For $m>1$ this contradicts
\[
 36m^2(9m^2+4)-13(4m^2+2)^2
                   =4(m^2-1)(29m^2+13)>0.
\]

For the exceptional case we need a little field arithmetic.
A cyclotomic field $\Q(\zeta_N)$ is generated by
$\zeta_N=e^{2\pi i/N}$. Its automorphisms send $\zeta_N$ to
$\zeta_N^a$ with $\gcd(a,N)=1$, preserving rational polynomial
relations. Averaging over the automorphisms that fix a subfield $K$
lands in $K$. For example, the two conjugates of $\zeta_3$ have
average $(\zeta_3+\zeta_3^2)/2=-1/2$.
The conductor of an abelian number field is the smallest $N$ for
which it lies in $\Q(\zeta_N)$. A root of order exactly $N$ is
primitive. A Ramanujan sum adds primitive roots without weights;
a quadratic Gauss sum weights roots by a quadratic character,
which modulo an odd prime is $1$ on nonzero squares, $-1$ on
nonsquares, and $0$ at zero. The weighted sums isolate the
square-root part of the average below.

The exceptional squareclass cannot simply be ignored. For a sum of
$N$ roots of unity equal to $-a-b\sqrt{10}$, with nonnegative
integers $a,b$, the conductor-forty orbit estimate gives
$N\geq a+8b$. Equality, when the squared-root sum also lies in
$\Q(\sqrt{10})$, forces that second sum to be $a$.

Here is the orbit mechanism behind this specialized arithmetic lemma.
Take a common cyclotomic field and form the union of all Galois images
of the root multiset under automorphisms fixing $\Q(\sqrt{10})$.
This gives an orbit-constant multiset without assuming that the
original multiplicities were constant on orbits. The imprimitive
quadratic Gauss-sum formulas separate two costs. Primitive orders
not divisible by forty contribute rational Ramanujan sums. Orders
divisible by forty have zero rational part, with nonzero quadratic
orbit sums $\pm\sqrt{10}$ costing at least eight roots each.
Thus a negative rational contribution of size $a$ costs at least
$a$ roots, independently of the $8b$ cost of the quadratic part.
Equality allows only copies of $-1$ and the negative quadratic orbit
of primitive fortieth roots. Squaring those eight roots runs once
through the primitive twentieth roots, whose sum is $\mu(20)=0$.
Here $\mu$ is the M\"obius function: it vanishes when its argument has
a squared prime factor, and otherwise is $(-1)^r$ for $r$ distinct
prime factors. Squaring $-1$ contributes one. Finally, Galois averaging
multiplies
the original second sum by the group order precisely because that
sum lies in $\Q(\sqrt{10})$. Dividing by that order gives $a$.

When $9m^2+4=10t^2$, reduction modulo two first makes $m$ even;
reduction modulo four then makes $t$ even. The conductor-forty
bound reads
\[
 4m^2+2\geq\frac{m^2}{2}+4mt.
\]
For $m\geq4$,
$t/m=\sqrt{(9+4/m^2)/10}>\sqrt{9/10}>15/16$, so the
right side is greater than $17m^2/4>4m^2+2$, a contradiction.
The only remaining positive possibility below four is $m=t=2$,
where the first-moment bound is an
equality and requires second sum $2$. But
\eqref{eq:c-indicator-second} gives $-2$ or $-2-4\sqrt{10}$.
Thus $m=1$.

\begin{proposition}[The resulting dimensions and twists]
The boundary fusion ring is the Haagerup ring, with
\[
 d=\frac{3+\sqrt{13}}2,\qquad s=1+d^2=2+3d.
\]
Its center has twelve simple objects, with quantum dimensions
\[
 1,\quad d^2,\quad s\text{ four times},\quad3d\text{ six times}.
\]
The six residual twists have exponents
$2,5,6,7,8,11$ modulo thirteen, divided by thirteen.
\end{proposition}

The constituent $X$ is the unique dimension-$d^2$ central
simple, later called $b$. Retain the name $Y$ for the original
embedded extension's dimension-$s$ constituent. The later
choice of a center equivalence and the canonical $\HH_3$
extension uses $a$ for that selected algebra's constituent.
The two twist-one dimension-$s$ simples, $a$ and $c_0$, have
the same unflavoured characters; the original $Y$ cannot be
identified with the selected $a$ from branching counts alone.
None of these central labels denotes the boundary invertible
object $\alpha$.

For $m=1$, the first moment is $(-1-\sqrt{13})/2$ and has exactly
six roots available. Average each root over a cyclotomic Galois
group fixing $\sqrt{13}$. The quadratic Gauss-sum estimate bounds
the absolute value of its $\sqrt{13}$ coefficient by $1/12$,
with equality only for primitive orders thirteen or twenty-six.
Thus the required coefficient $-1/2$ exhausts all six roots, each
at its negative extremum. This argument allows the fractional
orbit multiplicities that averaging may produce. The corresponding
six-root negative quadratic periods are, respectively,
\[
 \frac{-1-\sqrt{13}}2,\qquad
\frac{1-\sqrt{13}}2.
\]
Here the per-root estimate can be seen directly from the conductor.
For a primitive root of order $N$, the average over the subgroup
fixing $\sqrt{13}$ is rational unless $13$ divides $N$.
When it does, insert the quadratic character modulo thirteen
into the ordinary cyclotomic average. The untwisted numerator
is the Ramanujan sum $\mu(N)$, where $\mu$ is the M\"obius
function. The twisted numerator vanishes unless
$N=13m$, with $m$ squarefree and coprime to thirteen: a
redundant prime-power factor cancels by a geometric sum.
In the remaining case the Chinese remainder theorem reduces
the twisted numerator to $\pm\sqrt{13}$, using the quadratic
Gauss sum modulo thirteen. Division by
$\varphi(N)=12\varphi(m)$ gives an irrational coefficient of
absolute value $1/[12\varphi(m)]$. Equality with $1/12$
requires $m=1$ or $2$, giving exactly the two orders above.
Their rational coefficients are respectively $-1/12$ and
$1/12$. This explains both the six-root budget and the sign
test without assuming that the original multiset was itself
a union of Galois orbits.

The entire root budget is already needed for the irrational
coefficient, so there is no room for rational-only orbits or
cancellation. The prescribed rational part therefore selects order
thirteen. All original roots are among the six nonresidue roots.
These are linearly independent over $\Q$: a relation would give
a polynomial of degree at most twelve, with zero constant term,
divisible by $\Phi_{13}(t)=1+t+\cdots+t^{12}$, hence the zero
polynomial. Comparison with the quadratic period forces each root
to occur once. Since every simple was repeated $n_j^2$ times,
there are six such simples and each has $n_j=1$.
Multiplication by two
interchanges residues and nonresidues, so the second moment is
$(-1+\sqrt{13})/2$, giving $\varepsilon=1$.
This is why conductor thirteen appears in the center even though the
original lattice extension had order three.

\section{A ring is not yet a unitary tensor category}

The ring and twists determine the standard Haagerup modular-data target
in this setting, but modular data alone do not identify a braided
category in general. The later analytic argument proves complete
unitarity, allowing the extension algebra to be treated as a
commutative Q-system and its all-module category as a unitary fusion
category. The unitary classification of the untwisted six-object ring
then leaves two possibilities, $\HH_2$ and $\HH_3$.

These categories are Morita equivalent, so their centers are equivalent.
Consequently the chiral representation category is the Haagerup center
$Z(\HH_3)$. The original embedded $H/U$ all-module category is still
only identified as $\HH_2$ or $\HH_3$. To obtain an extension with
the specifically chosen $\HH_3$ boundary category, one transports its
canonical Lagrangian algebra through the center equivalence. The
resulting holomorphic VOA is again isomorphic to $V_{E_8}$, but the
embedding and its chosen boundary realization need not be the original
ones.

Keeping these last distinctions visible prevents a common shortcut:
``the numbers look Haagerup'' is not a tensor-category proof, and
``the center is Haagerup'' does not select a particular Morita
representative as the boundary category.

\section{Exercises}

\begin{exercise}\label{ex:c-branch-1}
Using $m=6$ and the neutral Fibonacci label, verify the top weight
$1/3$ for $x_4$. Compute the two summands separately.

\par\smallskip\noindent{\footnotesize\normalfont
\hyperref[sol:ex:c-branch-1]{Solution on page~\pageref*{sol:ex:c-branch-1}.}}\par
\end{exercise}
\begin{exercise}\label{ex:c-branch-2}
Check the coefficients of $x_0$, $x_5$, and $y_4$ in
\eqref{eq:c-top-certificate}. Explain why a positive bound for each
individual ratio $W/\dim$ would be weaker than using the current
character inequalities.

\par\smallskip\noindent{\footnotesize\normalfont
\hyperref[sol:ex:c-branch-2]{Solution on page~\pageref*{sol:ex:c-branch-2}.}}\par
\end{exercise}
\begin{exercise}\label{ex:c-branch-3}
If a semisimple object decomposes as $\bigoplus_i m_iX_i$ with
pairwise nonisomorphic simples, prove
$\dim\End(\bigoplus_i m_iX_i)=\sum_i m_i^2$.
Apply this to both branching alternatives for $H$.

\par\smallskip\noindent{\footnotesize\normalfont
\hyperref[sol:ex:c-branch-3]{Solution on page~\pageref*{sol:ex:c-branch-3}.}}\par
\end{exercise}
\begin{exercise}\label{ex:c-branch-4}
Show that the ring \eqref{eq:c-HI-m} has two one-dimensional
representations and one two-dimensional representation over $\C$.
Write the matrices in the two-dimensional representation and check
all defining relations.

\par\smallskip\noindent{\footnotesize\normalfont
\hyperref[sol:ex:c-branch-4]{Solution on page~\pageref*{sol:ex:c-branch-4}.}}\par
\end{exercise}
\begin{exercise}\label{ex:c-branch-5}
Verify the difference-of-squares factorization excluding $m>1$ outside
the exceptional squareclass. For $m=2$, compute both possible second
moments and compare them with the equality-case requirement.

\par\smallskip\noindent{\footnotesize\normalfont
\hyperref[sol:ex:c-branch-5]{Solution on page~\pageref*{sol:ex:c-branch-5}.}}\par
\end{exercise}
\begin{exercise}\label{ex:c-branch-6}
Use $d^2=1+3d$ to check that the squared dimensions of the twelve
central simples sum to $[3(1+d^2)]^2$. Distinguish this global
dimension from the total quantum dimension and from the rank twelve.

\par\smallskip\noindent{\footnotesize\normalfont
\hyperref[sol:ex:c-branch-6]{Solution on page~\pageref*{sol:ex:c-branch-6}.}}\par
\end{exercise}
